% LaTeX source for an arXiv-style conceptual paper on "Sunyata Mathematics (Ku-Mathematics)"
\documentclass[11pt]{article}
\usepackage{amsmath, amssymb, amsthm, amsfonts}
\usepackage{hyperref}
\usepackage{fullpage}
\usepackage{times}

\newcommand{\mweq}{\mathrel{\overset{\mathrm{mw}}{=}}}

\title{Toward a Theory of ``Sunyata Mathematics'':\\ A Foundational Framework for Observer-Dependent Abstraction}
\author{Masaomi Chiba\\ independent researcher}
\date{2025-12-05}

\begin{document}
\maketitle

\begin{abstract}
This paper introduces a conceptual framework called \emph{Sunyata Mathematics} (``Ku-Mathematics''),
intended as a pre-foundational layer underlying category theory, model theory, and other modern abstractions.
Sunyata Mathematics formalizes two key cognitive operations:
(i) the \emph{emptying} (``ku-ization'') of structures into observer-dependent coherent states, and
(ii) the \emph{coloring} (``shiki-ization'') of these states into determinate mathematical objects.
We argue that this bidirectional process provides a more primitive and flexible substrate than traditional foundations,
enabling a principled account of multi-universe reasoning, observer variability, and non-isomorphic translations.
We further propose that this framework is required to model certain extreme mathematical theories
that inherently rely on incompatible or non-translatable structures.
We define truth via dynamic stability rather than static equality and introduce the middle-way equivalence $\mweq$
through a convergence-based criterion tied to structural distortion $E(x)$.
\end{abstract}

\section{Introduction}
Classical foundations of mathematics---set theory, type theory, and category theory---presuppose that
mathematical structures exist in a single determinate universe equipped with stable logical rules.
However, several advanced theories (e.g., anabelian geometry and relative motives) expose phenomena
that cannot be internalized in a single universe without destructive loss of structure.

Recent LLM behavior further motivates this shift by exhibiting probabilistic, observer-dependent notions of truth,
suggesting that truth preservation itself may be dynamic rather than fixed.

We introduce \emph{Sunyata Mathematics} as a meta-level framework that explicitly models
(1) observer dependence, (2) multi-universe inconsistency, and
(3) non-resolution of certain structural relations.
This framework extends beyond category theory by incorporating non-translatable states and
systematically handling the alternation between undetermined and determined mathematical states.

The motivation also includes a shift from purely structural ``how''-descriptions toward an explicit account of
why certain identities should hold: static equality can mask structural mismatches across universes, dynamics,
or self-referential boundaries. Sunyata Mathematics treats such failures as non-middle fallacies (NMF) and
replaces rigid identity with convergence-based middle-way equivalence.

\section{Overview and Objectives}
Sunyata Mathematics is an optimization-oriented logical framework that treats truth as dynamic stability rather than static identity.
Its core objective is to reduce structural distortion $E(x)$ and resolve the non-middle fallacy by guiding structures toward middle-way stability.
Concretely, it aims to provide a principled account of multi-universe reasoning, to clarify the limits of direct translation between mathematical
universes, and to supply a formal basis for observer-variable inference in both human and machine reasoning.

Sunyata Mathematics also reframes truth as operational: truths are generated by the sunyata-ing/rupa-ing dynamics (as operational processes) rather than fixed once and for all.
This positions the framework as a pre-foundational layer that both generates mathematical universes and mediates translation between them,
each universe being stabilized around a distinct observer fixed point (Dfix0).
This amounts to a third truth regime distinct from deduction and induction, grounded in the outcomes of the operational cycle itself.

\section{Coherent and Decoherent Mathematical States}
We define two complementary modes:

\begin{itemize}
    \item \textbf{Coherent (Sunyata) States}: mathematical entities remain in an indeterminate state;
    morphisms, objects, and even universes are not yet fixed. These states represent ``pre-mathematical''
    abstraction.
    \item \textbf{Decoherent (Rupa) States}: specific mathematical structures become fixed;
    categories, sets, elements, and morphisms acquire determinate identity.
\end{itemize}

Sunyata Mathematics studies the controlled alternation between these states. Category theory corresponds to a fully decoherent phase, where composition, identity, and naturality are strictly enforced. In contrast, coherent states allow incompatible identifications and alternative observer positions.

\section{Middle-Way Equality and Structural Distortion}
Sunyata Mathematics defines truth not by a static equality but by dynamic stability, understood as convergence.
To make this explicit, we introduce a structural distortion function $E(x)$ that measures the deviation of a structure from the middle way.
The defining condition for middle-way stability is the equivalence
\[
E(x) = 0 \quad \Longleftrightarrow \quad x \mweq x_{\mathrm{mw}}.
\]
Here, $x_{\mathrm{mw}}$ is a middle-way state. The symbol $\mweq$ denotes middle-way equivalence.
We take $\mweq$ to be convergence-based: $x \mweq y$ holds when both converge to the same
middle-way state under admissible stabilization dynamics.

\section{Observer-Dependent Universes and Dfix0}
Let $F$ denote an observer configuration (``Dfix0''). A coherent state is defined relative to $D$.
Different observers may induce incompatible decoherent states even when starting from identical coherent seeds.

This formalizes the notion that mathematical truth, when abstract enough, becomes partially observer-relative.

\section{Non-Isomorphic Translations}
One motivation for Sunyata Mathematics is to give a principled account of non-isomorphic translations:
transitions between mathematical universes where no natural functor or equivalence exists.
Sunyata Mathematics explicitly allows ``colorings'' that cannot be compared within any category.

Structures where certain identifications are intentionally destroyed fit naturally into this framework.

\section{Sunyata Mathematics and Category Theory}
From the perspective of Sunyata Mathematics, category theory can be understood as a static abstraction of a partially emptying process.
It moves away from fixed essence by privileging morphisms and relations, but it does not internalize the observer (Dfix0) or the
dynamic alternation between emptying and coloring.
In this sense, category theory occupies an intermediate layer: it captures a ``Sunyata-leaning'' mathematical universe that is
not yet fully dynamical, while ordinary mathematics represents a further stage of fixation.
Sunyata Mathematics can therefore subsume category theory as a stabilized subsystem and provide an account of why such a
relation-centered universe is coherent within a broader, observer-aware framework.

Sunyata Mathematics also distinguishes local consistency from global coherence: a system can be internally consistent while
remaining misaligned with the broader dynamics that generate and stabilize its structures. This motivates treating coherence
as a first-class notion rather than assuming it is guaranteed by formal consistency alone.

\section{Toward Applications}
Potential applications include:

\begin{itemize}
    \item a meta-logic for multi-universe mathematics,
    \item a foundation for future AI systems capable of observer-variable reasoning,
    \item a cognitive model explaining creativity in high-level mathematics.
\end{itemize}

Beyond current LLMs, the framework suggests an AI paradigm that can (i) shift observer fixed points (Dfix0),
(ii) sustain both sunyata- and rupa-states without premature fixation, and (iii) traverse multiple mathematical universes
without forcing reductions to a single representation. These capabilities mark a distinct target for future AI design
grounded in the operational cycle of Sunyata Mathematics.

\section{Conclusion}
Sunyata Mathematics proposes a pre-foundational abstraction layer that models both coherent and decoherent phases of mathematical reasoning. Future work includes formalizing observer shifts, defining coherent-to-decoherent operators, and exploring connections with homotopy type theory and higher category theory.

\section*{References}
\begin{itemize}
    \item Masaomi Chiba, \emph{Dfix0 as a Middle-Way Fixed Point: An Operational Implementation of the Two Truths in Multi-Model Reasoning Systems}, 2025-12-01.
\end{itemize}

\end{document}
